% Options for packages loaded elsewhere
\PassOptionsToPackage{unicode}{hyperref}
\PassOptionsToPackage{hyphens}{url}
\documentclass[
  12pt,
]{article}
\usepackage{xcolor}
\usepackage{amsmath,amssymb}
\setcounter{secnumdepth}{-\maxdimen} % remove section numbering
\usepackage{iftex}
\ifPDFTeX
  \usepackage[T1]{fontenc}
  \usepackage[utf8]{inputenc}
  \usepackage{textcomp} % provide euro and other symbols
\else % if luatex or xetex
  \usepackage{unicode-math} % this also loads fontspec
  \defaultfontfeatures{Scale=MatchLowercase}
  \defaultfontfeatures[\rmfamily]{Ligatures=TeX,Scale=1}
\fi
\usepackage{lmodern}
\ifPDFTeX\else
  % xetex/luatex font selection
\fi
% Use upquote if available, for straight quotes in verbatim environments
\IfFileExists{upquote.sty}{\usepackage{upquote}}{}
\IfFileExists{microtype.sty}{% use microtype if available
  \usepackage[]{microtype}
  \UseMicrotypeSet[protrusion]{basicmath} % disable protrusion for tt fonts
}{}
\makeatletter
\@ifundefined{KOMAClassName}{% if non-KOMA class
  \IfFileExists{parskip.sty}{%
    \usepackage{parskip}
  }{% else
    \setlength{\parindent}{0pt}
    \setlength{\parskip}{6pt plus 2pt minus 1pt}}
}{% if KOMA class
  \KOMAoptions{parskip=half}}
\makeatother
\usepackage{color}
\usepackage{fancyvrb}
\newcommand{\VerbBar}{|}
\newcommand{\VERB}{\Verb[commandchars=\\\{\}]}
\DefineVerbatimEnvironment{Highlighting}{Verbatim}{commandchars=\\\{\}}
% Add ',fontsize=\small' for more characters per line
\newenvironment{Shaded}{}{}
\newcommand{\AlertTok}[1]{\textcolor[rgb]{1.00,0.00,0.00}{\textbf{#1}}}
\newcommand{\AnnotationTok}[1]{\textcolor[rgb]{0.38,0.63,0.69}{\textbf{\textit{#1}}}}
\newcommand{\AttributeTok}[1]{\textcolor[rgb]{0.49,0.56,0.16}{#1}}
\newcommand{\BaseNTok}[1]{\textcolor[rgb]{0.25,0.63,0.44}{#1}}
\newcommand{\BuiltInTok}[1]{\textcolor[rgb]{0.00,0.50,0.00}{#1}}
\newcommand{\CharTok}[1]{\textcolor[rgb]{0.25,0.44,0.63}{#1}}
\newcommand{\CommentTok}[1]{\textcolor[rgb]{0.38,0.63,0.69}{\textit{#1}}}
\newcommand{\CommentVarTok}[1]{\textcolor[rgb]{0.38,0.63,0.69}{\textbf{\textit{#1}}}}
\newcommand{\ConstantTok}[1]{\textcolor[rgb]{0.53,0.00,0.00}{#1}}
\newcommand{\ControlFlowTok}[1]{\textcolor[rgb]{0.00,0.44,0.13}{\textbf{#1}}}
\newcommand{\DataTypeTok}[1]{\textcolor[rgb]{0.56,0.13,0.00}{#1}}
\newcommand{\DecValTok}[1]{\textcolor[rgb]{0.25,0.63,0.44}{#1}}
\newcommand{\DocumentationTok}[1]{\textcolor[rgb]{0.73,0.13,0.13}{\textit{#1}}}
\newcommand{\ErrorTok}[1]{\textcolor[rgb]{1.00,0.00,0.00}{\textbf{#1}}}
\newcommand{\ExtensionTok}[1]{#1}
\newcommand{\FloatTok}[1]{\textcolor[rgb]{0.25,0.63,0.44}{#1}}
\newcommand{\FunctionTok}[1]{\textcolor[rgb]{0.02,0.16,0.49}{#1}}
\newcommand{\ImportTok}[1]{\textcolor[rgb]{0.00,0.50,0.00}{\textbf{#1}}}
\newcommand{\InformationTok}[1]{\textcolor[rgb]{0.38,0.63,0.69}{\textbf{\textit{#1}}}}
\newcommand{\KeywordTok}[1]{\textcolor[rgb]{0.00,0.44,0.13}{\textbf{#1}}}
\newcommand{\NormalTok}[1]{#1}
\newcommand{\OperatorTok}[1]{\textcolor[rgb]{0.40,0.40,0.40}{#1}}
\newcommand{\OtherTok}[1]{\textcolor[rgb]{0.00,0.44,0.13}{#1}}
\newcommand{\PreprocessorTok}[1]{\textcolor[rgb]{0.74,0.48,0.00}{#1}}
\newcommand{\RegionMarkerTok}[1]{#1}
\newcommand{\SpecialCharTok}[1]{\textcolor[rgb]{0.25,0.44,0.63}{#1}}
\newcommand{\SpecialStringTok}[1]{\textcolor[rgb]{0.73,0.40,0.53}{#1}}
\newcommand{\StringTok}[1]{\textcolor[rgb]{0.25,0.44,0.63}{#1}}
\newcommand{\VariableTok}[1]{\textcolor[rgb]{0.10,0.09,0.49}{#1}}
\newcommand{\VerbatimStringTok}[1]{\textcolor[rgb]{0.25,0.44,0.63}{#1}}
\newcommand{\WarningTok}[1]{\textcolor[rgb]{0.38,0.63,0.69}{\textbf{\textit{#1}}}}
\usepackage{longtable,booktabs,array}
\usepackage{calc} % for calculating minipage widths
% Correct order of tables after \paragraph or \subparagraph
\usepackage{etoolbox}
\makeatletter
\patchcmd\longtable{\par}{\if@noskipsec\mbox{}\fi\par}{}{}
\makeatother
% Allow footnotes in longtable head/foot
\IfFileExists{footnotehyper.sty}{\usepackage{footnotehyper}}{\usepackage{footnote}}
\makesavenoteenv{longtable}
\setlength{\emergencystretch}{3em} % prevent overfull lines
\providecommand{\tightlist}{%
  \setlength{\itemsep}{0pt}\setlength{\parskip}{0pt}}
\usepackage{bookmark}
\IfFileExists{xurl.sty}{\usepackage{xurl}}{} % add URL line breaks if available
\urlstyle{same}
\hypersetup{
  hidelinks,
  pdfcreator={LaTeX via pandoc}}

\author{SCX}
\date{2026}

\begin{document}

\section{Spring Self-Evolution Algorithm ---
Mathematical Genealogy}\label{spring-ux81eaux8fdbux5316ux7b97ux6cd5-ux6570ux5b66ux8c31ux7cfbux5b66}

\begin{quote}
\textbf{Spring} = Spring Self-Evolving Gatekeeper (SSEG)
\textbf{Status}: Paper 2 (Nature Computational Science target) core algorithm
\textbf{Document Type}: Theoretical Genealogy ---
Mathematical domain classification, key idea origins, historical development timeline, and comparison with prior work
\textbf{Date}: 2026-06-28
\end{quote}

\begin{center}\rule{0.5\linewidth}{0.5pt}\end{center}

\subsection{Table of Contents}\label{ux76eeux5f55}

\begin{enumerate}
\def\labelenumi{\arabic{enumi}.}
\tightlist
\item
  \hyperref[1-ux6570ux5b66ux9886ux57dfux5f52ux5c5e]{Mathematical Domain Classification}
\item
  \hyperref[2-ux5173ux952eux601dux60f3ux4e0eux8bc1ux660eux5de5ux5177]{Key Ideas and Proof Tools}
\item
  \hyperref[3-ux5386ux53f2ux53d1ux5c55ux65f6ux95f4ux7ebf]{Historical Development Timeline}
\item
  \hyperref[4-spring-ux4e0eux524dux4ebaux7684ux5bf9ux6bd4ux8868]{Spring
  Comparison with Prior Work}
\item
  \hyperref[5-spring-ux72ecux6709ux7684ux8d21ux732e]{Spring's Unique Contributions}
\end{enumerate}

\begin{center}\rule{0.5\linewidth}{0.5pt}\end{center}

\subsection{1. Mathematical Domain Classification}\label{ux6570ux5b66ux9886ux57dfux5f52ux5c5e}

Spring
The Spring self-evolution algorithm does not belong to a single mathematical domain, but spans seven classical mathematical domains, synthesizing a completely new structure at their intersection.

\subsubsection{1.1 Stochastic Approximation (Stochastic
Approximation)}\label{ux968fux673aux903cux8fd1-stochastic-approximation}

\textbf{Origin}: Robbins \& Monro (1951). ``A stochastic approximation
method.'' \emph{Annals of Mathematical Statistics}, 22(3), 400-407.

\textbf{Role in Spring}: NEP student parameter \(\theta_t\) 'sonline update follows
Robbins-Monro recursion:

\[\theta_{t+1} = \theta_t - \alpha_t \nabla_\theta \ell(f_{\theta_t}(x_t), y_t)\]

where the learning rate \(\alpha_t\) satisfies the standard RM
conditions:\(\sum \alpha_t = \infty\)，\(\sum \alpha_t^2 < \infty\)。

Spring generalizes the standard RM to\textbf{coupled two-timescale}setting： -
\textbf{Fast timescale}: student  \(\theta_t\) at 固定 Gatekeeper \(S_t\)
under with Rate \(\alpha_t\) Update - \textbf{Slow timescale}: Gatekeeper \(S_t\)
under the student quasi-steady-state at rate \(\beta_t = o(\alpha_t)\) Update

This corresponds to Borkar (2008)'s two-timescale stochastic approximation framework, where condition C6'
(\(\beta_t = o(\alpha_t)\))
is the key innovation---it simultaneously guarantees local convergence of the student and bounded cumulative distribution shift.

\textbf{References}: - Robbins \& Monro (1951) - Kushner \& Yin (2003).
\emph{Stochastic Approximation and Recursive Algorithms and
Applications} - Borkar (2008). \emph{Stochastic Approximation: A
Dynamical Systems Viewpoint} - Polyak \& Juditsky (1992). ``Acceleration
of stochastic approximation by averaging''

\subsubsection{1.2 Dynamical Systems / Lyapunov
Stability}\label{ux52a8ux529bux7cfbux7edf-lyapunov-ux7a33ux5b9aux6027}

\textbf{Origin}: Lyapunov, A. M. (1892). \emph{The General Problem of the
Stability of Motion}. Kharkov Mathematical Society.

\textbf{Role in Spring}: Spring self-evolution loopis adiscrete dynamical system：

\[z_{t+1} = \Phi(z_t), \quad z_t = (S_t, M_t, \theta_t) \in \mathcal{Z}\]

where \(\Phi = (\Phi_S, \Phi_M, \Phi_\theta)\)
Yesupdate operator，\(\mathcal{Z} = \mathcal{F} \times \mathcal{M} \times \Theta\)
Yesevolution state space。

Lyapunov Function候选as （On a fixed reference set \(M_0\) on evaluation):

\[\Psi(S_t, \theta_t) = \frac{1}{|M_0|} \sum_{x \in M_0} (S_t(x, y(x)) - \hat{C}(x))^2 + \lambda \cdot \frac{1}{|V_0|} \sum_{(x,y) \in V_0} \ell(f_{\theta_t}(x), y)\]

Lyapunov under 降性质
\(\mathbb{E}[\Psi_{t+1} \mid \mathcal{F}_t] \leq \Psi_t\) via  Theorem
12.5 at 参考集重放Mechanismunder is CompleteProof。non-动点 \((S^*, \theta^*)\)
满足自洽耦合方程，\(\omega\)-极限集as 单点。

Spring
make using 了动力System'sCompletetoolchain: attractors, basins, phase diagram (four convergencemechanisms), \(\omega\)-limit set.

\textbf{References}: - Lyapunov (1892). \emph{The General Problem of the
Stability of Motion} - Khalil (2002). \emph{Nonlinear Systems} (3rd ed.)
- Strogatz (2018). \emph{Nonlinear Dynamics and Chaos} (2nd ed.) -
Conley (1978). \emph{Isolated Invariant Sets and the Morse Index}

\subsubsection{1.3 贝叶斯推断 /
Martingale Theory}\label{ux8d1dux53f6ux65afux63a8ux65ad-ux9785ux7406ux8bba}

\textbf{Origin}: Doob, J. L. (1953). \emph{Stochastic Processes}. Wiley.

\textbf{Role in Spring}: Gatekeeper scoring function \(S_t\)
'sevolutionis 解释as 贝叶斯后验all 值Update：

\[S_t(x,y) = \mathbb{E}_{S \sim P_t}[S(x,y)], \quad P_{t+1}(S) \propto P_t(S) \cdot \mathcal{L}(\mathcal{D}_{t+1} \mid S)\]

定理 (Doob Martingale Convergence): 序列 \(\{S_t\}_{t \geq 0}\) Yes关于滤子
\(\mathcal{F}_t = \sigma(\mathcal{M}_1, \dots, \mathcal{M}_t)\) 's
\([0,1]\)-has 界鞅。by  Doob Martingale Convergence定理：

\[S_t(x,y) \xrightarrow{a.s.} S_\infty(x,y), \quad \forall (x,y)\]

重must 限制：Martingale Propertymust 求model is correct（\(S^*\) in the prior \(P_0\)
's支撑内）。at model misspecification设定under ，后验Convergenceto  KL 投影 \(S^\dagger\) rather than the true value
\(S^*\)。Spring honestly acknowledges this limitation.

\textbf{References}: - Doob (1953). \emph{Stochastic Processes} - Doob
(1949). ``Application of the theory of martingales'' - Ghosh \&
Ramamoorthi (2003). \emph{Bayesian Nonparametrics} - van der Vaart
(1998). \emph{Asymptotic Statistics} - Kleijn \& van der Vaart (2012).
``The Bernstein-von Mises theorem under misspecification''

\subsubsection{1.4 at 线Learning /
Regret Analysis}\label{ux5728ux7ebfux5b66ux4e60-ux9057ux61beux5206ux6790}

\textbf{Origin}: Zinkevich, M. (2003). ``Online convex programming and
generalized infinitesimal gradient ascent.'' \emph{ICML 2003}.

\textbf{Role in Spring}: Gatekeeper 'sat 线梯度under 降 (OGD)
update yields标准regret bound：

\[\text{Regret}_T \leq 2GR\sqrt{T} \quad \text{(No Delay)}\]

Spring generalizes this bound to the\textbf{delayed feedback}setting (the NEP student's feedback delay
\(d_t \leq D_{\max}\)):

\[\text{Regret}_T^{\text{delay}} \leq 2GR\sqrt{T} + G^2 D_{\max} \sqrt{T}\]

Further generalized to the\textbf{growing memory bank}setting------memory bank \(M_t\)
's单调增长提供has 效sample complexity改进：

\[\text{Regret}_T^{\text{eff}} \leq \text{Regret}_T^{\text{delay}} + O\left(T \sqrt{\frac{\text{VC}(\mathcal{H})}{N_T}}\right)\]

As
\(N_T \to \infty\)，System满足\textbf{Asymptotic No-Regret}：\(\frac{1}{T}\text{Regret}_T^{\text{eff}} \to 0\)。

\textbf{References}: - Zinkevich (2003). ``Online convex programming'' -
Hazan (2016). ``Introduction to online convex optimization'' - Joulani
et al.~(2016). ``Online learning under delayed feedback'' -
Shalev-Shwartz (2012). ``Online learning and online convex
optimization'' - Cesa-Bianchi \& Lugosi (2006). \emph{Prediction,
Learning, and Games}

\subsubsection{1.5 concentration inequalities}\label{ux96c6ux4e2dux4e0dux7b49ux5f0f}

\textbf{Origin}: - Hoeffding, W. (1963). ``Probability inequalities for
sums of bounded random variables.'' \emph{JASA}, 58(301), 13-30. -
Chernoff, H. (1952). ``A measure of asymptotic efficiency for tests of a
hypothesis.'' \emph{Annals of Mathematical Statistics}, 23(4), 493-507.

\textbf{Role in Spring}: at  Spring 's多个critical locations：

\begin{enumerate}
\def\labelenumi{\arabic{enumi}.}
\item
  \textbf{Concentration of the Consensus Score}: Expert consensus
  \(\hat{C}(x) = \frac{1}{M}\sum_{m=1}^M \mathbf{1}\{\ell(f_m(x), y) > \tau\}\)
  Hoeffding tail bound:
  \[\mathbb{P}(|C - \mathbb{E}[C]| > t) \leq 2\exp(-2M t^2)\] When A2
  is violated due to expert correlation,\(M\) is replaced by the effective number of experts
  \(M_{\text{eff}} = M/(1 + (M-1)\bar{\rho})\)。
\item
  \textbf{Chernoff 信息}: detection problem'soptimal exponentialRateby  Chernoff 信息
  \(\kappa = \text{KL}(\theta^*\|p_0) = \text{KL}(\theta^*\|p_1)\)
  , applied to Theorem 4'.
\item
  \textbf{Bahadur-Rao Exact Asymptotics}: Bernoulli lattice correctionbecause 子
  \((1-e^{-\lambda^*})^{-1}\) replaces the standard Chernoff prefactor, tightening Theorem 4'
  's exact constant by 5--12\%.
\end{enumerate}

\textbf{References}: - Hoeffding (1963) - Chernoff (1952) - Bahadur \& Rao
(1960). ``On deviations of the sample mean'' - Dembo \& Zeitouni (2010).
\emph{Large Deviations Techniques and Applications}

\subsubsection{1.6 Information Theory}\label{ux4fe1ux606fux8bba}

\textbf{Origin}: - Fano, R. M. (1961). \emph{Transmission of Information:
A Statistical Theory of Communications}. MIT Press. - Cover, T. M. \&
Thomas, J. A. (2006). \emph{Elements of Information Theory} (2nd ed.).
Wiley.

\textbf{Role in Spring}:

\begin{enumerate}
\def\labelenumi{\arabic{enumi}.}
\item
  \textbf{Fano non-etc.式}: using 于 Theorem 2（Weak Feature Failure）'slower bound------when feature
  \(\phi\) and true state \(S\) have mutual information \(I(\phi; S) \leq \delta\)
  时，任何state estimationAlgorithmhas ：
  \[P(\hat{S} \neq S) \geq \frac{H(S) - \delta - \log 2}{\log K}\]
\item
  \textbf{Data Processingnon-etc.式 (DPI)}: using 于 Theorem 2 Proof Step
  3------the Markov chain \(\phi(X) \to \hat{S} \to \hat{z}\) 意味着 TV
  does not increase.
\item
  \textbf{Pinsker non-etc.式}: will mutual information bound \(\delta = I(\phi; S)\)
  to a total variation bound \(\text{TV} \leq \sqrt{\delta/2}\)。
\item
  \textbf{KL 收缩}: at correct modelunder ，后验Convergence
  \(D_{\text{KL}}(P^* \| P_t) \xrightarrow{a.s.} 0\)。
\end{enumerate}

\textbf{References}: - Fano (1961) - Cover \& Thomas (2006) - Pinsker
(1964). \emph{Information and Information Stability of Random Variables
and Processes}

\subsubsection{1.7 Minimax Theory}\label{minimax-ux7406ux8bba}

\textbf{Origin}: Le Cam, L. (1973). ``Convergence of estimates under
dimensionality restrictions.'' \emph{Annals of Statistics}, 1(1), 38-53.

\textbf{Role in Spring}:

\begin{enumerate}
\def\labelenumi{\arabic{enumi}.}
\item
  \textbf{Le Cam two-point method}: 构造两个难with 区分'sDistribution
  \(P_0 = \text{Bern}(\mu_s)^{\otimes M}\)（Clean）和
  \(P_1 = \text{Bern}(1 - C_{\text{bal}}\mu_s/(K-1))^{\otimes M}\)（Noise），利using 
  Hellinger tensorization 获得 minimax lower bound。
\item
  \textbf{Neyman-Pearson 归约}: 通过 Neyman-Pearson 引理will ArbitraryAlgorithm归约as 
  Bayes 检验，give 出 Theorem 4' 'sOptimal性。
\item
  \textbf{Hellinger 精确张量化}: at 
  A2（Conditional Independence）under ，\(H^2(P_0^{\otimes M}, P_1^{\otimes M}) = 1 - (1 - H^2(P_0, P_1))^M\)------精确etc.式，Nonon-etc.式松弛。
\end{enumerate}

\textbf{References}: - Le Cam (1973) - Le Cam (1986). \emph{Asymptotic
Methods in Statistical Decision Theory} - Tsybakov (2009).
\emph{Introduction to Nonparametric Estimation}

\begin{center}\rule{0.5\linewidth}{0.5pt}\end{center}

\subsection{2.
Key Ideas and Proof Tools}\label{ux5173ux952eux601dux60f3ux4e0eux8bc1ux660eux5de5ux5177}

\subsubsection{\texorpdfstring{2.1 Two-Timescale Separation
(\(\beta_t = o(\alpha_t)\))}{2.1 Two-Timescale Separation (\textbackslash beta\_t = o(\textbackslash alpha\_t))}}\label{ux53ccux65f6ux95f4ux5c3aux5ea6ux5206ux79bb-beta_t-oalpha_t}

Spring 最关键'sarchitectural innovation. Condition C6' must 求 Gatekeeper update rate \(\beta_t\)
Yesstudent update rate \(\alpha_t\) 'slower-order infinitesimal:

\[\beta_t = o(\alpha_t), \quad \text{i.e.} \quad \lim_{t\to\infty} \frac{\beta_t}{\alpha_t} = 0\]

\textbf{目's}: Between any two Gatekeeper
updates，student performs asymptotically infinitely many gradient steps，at before distribution shift, convergingto when 前Distribution's局部最小值。

\textbf{Canonical Schedule}: \(\alpha_t = t^{-a}\)，\(\beta_t = t^{-b}\)，where
\(\frac{1}{2} < a < 1 < b\)。例such as  \(a = 0.6\)，\(b = 1.2\)： -
\(\sum \alpha_t = \infty\)（student exploration），\(\sum \alpha_t^2 < \infty\)（noise control）
-
\(\sum \beta_t < \infty\)（has 限总Distribution Shift），\(\sum \beta_t^2 < \infty\)（variance control）
- \(\beta_t / \alpha_t = t^{-0.6} \to 0\)（Timescale Separation）

\subsubsection{2.2
参考集重放打破Selection Bias}\label{ux53c2ux8003ux96c6ux91cdux653eux6253ux7834ux9009ux62e9ux504fux5dee}

Spring 最深刻'sProof工具。Theorem 12.2
Proof：\textbf{Without reference set replay, Lyapunov
under 降at 渐近意义under Yesnon-can be can 's}------Gatekeeper
can reduce its apparent loss by only accepting''easy''samples to reduce apparent loss，while non-真正改进判别can 力。

Resolution方案 (Theorem 12.5, CompleteProof): 1. \textbf{Student Side}: importance sampling weights
\(w(x) = P_0(x) / P_{S_t}(x)\)，make 得has 效training distributionas reference distribution \(P_0\) 2.
\textbf{Gatekeeper Side}: On a fixed reference set \(M_0\) on Computation
SCXUpdate，eliminating alignment bias 3. \textbf{Cost}: Importance weight variance inflation factor
\(W \leq 1/\varepsilon_{\text{explore}} < \infty\)

\subsubsection{2.3 Lipschitz-Almost-Everywhere
Lemma}\label{lipschitz-almost-everywhere-ux5f15ux7406}

Resolves a fundamental mathematical obstacle (Lemma SE-1.0, B9 fix):

\textbf{Problem}: consensus score
\(C(x) = \frac{1}{M}\sum_{m=1}^M \mathbf{1}\{\ell(f_m(x), y) > \tau\}\)
at the threshold \(\tau\) 处non-连续（指示Function跳跃）。

\textbf{Resolution}: at conditions C5（conditions i.i.d.
sampling), the discontinuity set has probability measure zero, because:
\[\mathbb{P}(\exists m: \ell(f_m(X), Y) = \tau) \leq \sum_{m=1}^M \mathbb{P}(\ell(f_m(X), Y) = \tau) = 0\]
for 于连续Distributionon 's光滑Function
\(\ell(f_m(x), y)\)。期望non-受零测集修改's影响，Lyapunov Analysis保留has 效。

\subsubsection{2.4 贝叶斯后验Update作as  Gatekeeper
Evolution}\label{ux8d1dux53f6ux65afux540eux9a8cux66f4ux65b0ux4f5cux4e3a-Gatekeeper-ux8fdbux5316}

Gatekeeper 'sEvolutionis 解释as 序贯Bayesian update：

\[P_{t+1}(S) \propto P_t(S) \cdot \prod_{(x,y) \in \mathcal{D}_{t+1}} S(x,y)^{a} (1 - S(x,y))^{1-a}\]

this 产生了两个强大'stheoretical properties： 1. \textbf{Martingale Property}:
\(\mathbb{E}[S_{t+1} \mid \mathcal{F}_t] = S_t\)（Tower Property），by  Doob
theorem guarantees a.s. Convergence 2. \textbf{KL 收缩}:
at correct modelunder ，\(D_{\text{KL}}(P^* \| P_t) \to 0\)

\subsubsection{2.5 Lyapunov
Descent with Importance Sampling}\label{lyapunov-ux4e0bux964dux4e0eux91cdux8981ux6027ux91c7ux6837}

Theorem 12.5 'scomplete proof structure establishes
\(\mathbb{E}[\Delta\Psi_t \mid \mathcal{F}_t] \leq 0\)：

\[\mathbb{E}[\Delta\Psi_t \mid \mathcal{F}_t] \leq -\alpha_t \|\nabla L_0(\theta_t)\|^2 - \beta_t \cdot 2\rho_{\text{ideal}} \cdot \|\Delta_t^{\text{ideal}}\|_{M_0} \cdot \|S_t - \hat{C}\|_{M_0} + O(\alpha_t^2 W + \beta_t^2)\]

Decomposition into four terms: - \textbf{(A) Student Improvement}:
\(\Delta_{\text{student}}\), eliminating via importance sampling \(D_{\text{static}}\)
gap (\textbf{already Proof}) - \textbf{(B) Gatekeeper Improvement}:
\(\Delta_{\text{Gatekeeper}}\), by computing \(M_0\) on Computation SCXUpdate
eliminating alignment bias（\textbf{already Proof}) - \textbf{(C) Distribution Shift}:
\(\Delta_{\text{selection}} = O(\beta_t) = o(\alpha_t)\), controlled by two-timescale (\textbf{already Proof}）
- \textbf{(D) Cross-Coupling}:
\(\Delta_{\text{cross}} = 0\), due to reference set separability (\textbf{already Proof}）

\subsubsection{\texorpdfstring{2.6 \(O(t^{-a})\) Convergence率and  Polyak
平all }{2.6 O(t\^{}\{-a\}) Convergence率and  Polyak 平all }}\label{ot-a-ux6536ux655bux7387ux4e0e-polyak-ux5e73ux5747}

at 强凸conditionsunder ，student 'sConvergence率as  \(O(t^{-a})\)，where \(a \in (0.5, 1]\)
as learning rate exponent。

\textbf{Polyak-Ruppert Averaging}
(\(\bar{\theta}_t = \frac{1}{t}\sum_{i=1}^t \theta_i\)) will Rate提升to 
\(O(t^{-1})\)------and  minimax lower bound \(\Omega(t^{-1})\)
matching'sOptimalRate------for any \(a \in (0.5, 1)\) and Norequires 精确调谐。

\subsubsection{2.7
Characterization of Four Failure Modes}\label{ux56dbux79cdux5931ux6548ux6a21ux5f0fux523bux753b}

Spring systematically characterizesfour canonical failure modes：

\begin{longtable}[]{@{}
  >{\raggedright\arraybackslash}p{(\linewidth - 4\tabcolsep) * \real{0.3000}}
  >{\raggedright\arraybackslash}p{(\linewidth - 4\tabcolsep) * \real{0.3000}}
  >{\raggedright\arraybackslash}p{(\linewidth - 4\tabcolsep) * \real{0.4000}}@{}}
\toprule\noalign{}
\begin{minipage}[b]{\linewidth}\raggedright
Mode
\end{minipage} & \begin{minipage}[b]{\linewidth}\raggedright
Mechanism
\end{minipage} & \begin{minipage}[b]{\linewidth}\raggedright
Severity
\end{minipage} \\
\midrule\noalign{}
\endhead
\bottomrule\noalign{}
\endlastfoot
\textbf{1. Premature Freezing} & \(\beta_t\) decays too quickly before the student provides corrective signals &
in etc.（双Accuracyunder 罕见） \\
\textbf{2. Backlog} & Data arrivalRate超过 Gatekeeper 评分吞吐量 &
高（大datasetunder 实际存at ） \\
\textbf{3. Client Divergence} & 多个non-动点 + non-同initial conditions/data streams →
irreconcilable divergence & Medium (multi-client deployment only) \\
\textbf{4. Adversarial Poisoning} & for 手注入精心构造'sSample；self-evolution loop放大污染 &
High (severe with amplification) \\
\end{longtable}

each 种Mode配has 形式化conditions、退化Rate和缓解策略。

\begin{center}\rule{0.5\linewidth}{0.5pt}\end{center}

\subsection{3.
Historical Development Timeline}\label{ux5386ux53f2ux53d1ux5c55ux65f6ux95f4ux7ebf}

\begin{verbatim}
1892 ─ Lyapunov StabilityTheory
      A. M. Lyapunov, The General Problem of the Stability of Motion
      动力System定性Analysis's基础——non-动点、吸引子、Lyapunov Function
      ↓ Spring make using : Ψ(S_t, θ_t) as  Lyapunov Function；non-动点Analysis；吸引子盆地
      
1933 ─ Kolmogorov Strong Law of Large Numbers
      A. N. Kolmogorov, Grundbegriffe der Wahrscheinlichkeitsrechnung
      独立同Distribution序列's几乎必然Convergence基础
      ↓ Spring make using : consensus statistic C(x) 'slarge-samplelinesas 
      
1951 ─ Robbins-Monro Stochastic Approximation
      H. Robbins & S. Monro, "A stochastic approximation method"
      Stochastic Gradient Descent's数学基础；∑α_t=∞, ∑α_t²<∞ conditions
      ↓ Spring make using : θ_t update follows RM 递归；双时间尺度extension
      
1952 ─ Chernoff 信息
      H. Chernoff, "A measure of asymptotic efficiency for tests"
      Assumptions检验'soptimal exponentialRate
      ↓ Spring make using : Theorem 4' Detection's精确指数Rate

1953 ─ Doob Martingale Convergence
      J. L. Doob, Stochastic Processes
      has 界鞅's几乎必然Convergence
      ↓ Spring Used for: S_t as a [0,1]-bounded martingale → S_t → S_∞ a.s.

1960 ─ Bahadur-Rao Exact Asymptotics
      R. R. Bahadur & R. R. Rao, "On deviations of the sample mean"
      Sampleall 值尾概率's精确二阶渐近（含lattice correction）
      ↓ Spring Used for: exact constant for Theorem 4'; Bernoulli lattice correction
      
1961 ─ Fano non-etc.式
      R. M. Fano, Transmission of Information
      Information Theorylower bound——互信息and error probability's关系
      ↓ Spring make using : Theorem 2 弱featurelower bound
      
1963 ─ Hoeffding non-etc.式
      W. Hoeffding, "Probability inequalities for sums of bounded random variables"
      Exponential concentration for sums of bounded independent variables
      ↓ Spring Used for: concentration bound for consensus score; main tool for Theorem 1
      
1964 ─ Pinsker non-etc.式
      M. S. Pinsker, Information and Information Stability
      TV ≤ √(KL/2) --- bridge between information divergence and statistical distance
      ↓ Spring Used for: δ → TV conversion in Theorem 2
      
1973 ─ Le Cam Minimax
      L. Le Cam, "Convergence of estimates under dimensionality restrictions"
      two-point methodlower bound、convex hull argument、minimax optimality framework
      ↓ Spring make using : Theorem 4' lower bound；Hellinger tensorization
      
1981 ─ Pollard k-means 一致性
      D. Pollard, "Strong consistency of k-means clustering"
      Almost sure consistency of k-means
      ↓ Spring make using : Theorem 5 state discovery'sTheory基础

1985 ─ QWERTY Lock-In (David)
      P. A. David, "Clio and the Economics of QWERTY"
      Pathdependenceand 次优技术Lock-In'svia 济学Theory
      ↓ Conceptual analogy: Spring 's多non-动点Structure——Gatekeeper can be can Lock-In于次优自洽配置
      
1989 ─ Increasing Returns (Arthur)
      W. B. Arthur, "Competing technologies, increasing returns, and lock-in by historical events"
      Multiple equilibria in positive feedback loops
      ↓ Conceptual analogy: Spring Gatekeeper 'sSelection Bias循环——initial bias is amplified throughSelf-Evolutionis 放大

2003 ─ Online Convex Programming (Zinkevich)
      M. Zinkevich, "Online convex programming and generalized infinitesimal gradient ascent"
      Regret bound O(√T) for online gradient descent
      ↓ Spring make using : Gatekeeper OGD update、delayed feedbackextension、growing memory bankextension

2010 ─ Domain自适should Theory
      Ben-David et al., "A theory of learning from different domains"
      Trainingand TestDistributionnon-matching's泛化界
      ↓ Spring make using : D_static 间隙'sDomain自适should 界 (Lemma 12.1)

2017 ─ AlphaZero Self-Play
      D. Silver et al., "Mastering Chess and Shogi by Self-Play..."
      Achieves superhuman level through self-play, policy iteration, Monte Carlo tree search
      ↓ Spring 类似: self-evolution loopStructure、monotone improvement guarantee、memory replay
      ↓ Spring non-同: 处理真实physicalData（non-合成）、not yet 知环境动态、formal convergence proof

2026 ─ Spring Self-Evolution Gatekeeper (This Work)
      SCX Research Group, "Spring: A Self-Evolving Gatekeeper with Provable Convergence"
      Self-Evolutionevaluation criteria（non-only 自改进策略）、monotonically growing memory、
      resurrection mechanism——"Winter does not kill --- it waits for spring"
\end{verbatim}

\begin{center}\rule{0.5\linewidth}{0.5pt}\end{center}

\subsection{4. Spring
Comparison with Prior Work}\label{spring-ux4e0eux524dux4ebaux7684ux5bf9ux6bd4ux8868}

\subsubsection{4.1 Core Comparison}\label{ux6838ux5fc3ux5bf9ux6bd4}

\begin{longtable}[]{@{}
  >{\raggedright\arraybackslash}p{(\linewidth - 10\tabcolsep) * \real{0.2500}}
  >{\raggedright\arraybackslash}p{(\linewidth - 10\tabcolsep) * \real{0.1154}}
  >{\raggedright\arraybackslash}p{(\linewidth - 10\tabcolsep) * \real{0.1731}}
  >{\raggedright\arraybackslash}p{(\linewidth - 10\tabcolsep) * \real{0.2115}}
  >{\raggedright\arraybackslash}p{(\linewidth - 10\tabcolsep) * \real{0.1346}}
  >{\raggedright\arraybackslash}p{(\linewidth - 10\tabcolsep) * \real{0.1154}}@{}}
\toprule\noalign{}
\begin{minipage}[b]{\linewidth}\raggedright
Algorithm / Framework
\end{minipage} & \begin{minipage}[b]{\linewidth}\raggedright
Year
\end{minipage} & \begin{minipage}[b]{\linewidth}\raggedright
Self-Evolving?
\end{minipage} & \begin{minipage}[b]{\linewidth}\raggedright
Convergence Proof?
\end{minipage} & \begin{minipage}[b]{\linewidth}\raggedright
Memory?
\end{minipage} & \begin{minipage}[b]{\linewidth}\raggedright
Domain
\end{minipage} \\
\midrule\noalign{}
\endhead
\bottomrule\noalign{}
\endlastfoot
\textbf{Robbins-Monro} & 1951 & No & Yes（渐近 a.s.） & No &
通using Stochastic Approximation \\
\textbf{Hedge / 专家建议} & 1997 & No & Yes（遗憾 O(√T)） & No &
at 线决策 \\
\textbf{主动Learning} & 2009 & No（Fixed Query Strategy） & Yes（Label Complexity） &
Label pool & 分类（单Model） \\
\textbf{Bayesian Optimization} & 2016 & No（Fixed Prior） & Yes（遗憾） & GP posterior &
Black-box optimization \\
\textbf{AlphaZero} & 2017 & Yes（Self-Play） & No（via 验性）* &
Replay buffer (with replacement) & Games (known environment) \\
\textbf{Solomonoff 归纳} & 1964 & N/A（Theory极限） & Yes（for can be Computation环境） &
All history & 通using 归纳（non-can be Computation） \\
\textbf{Spring (Cercis)} & 2026 & \textbf{Yes}（Gatekeeper co-evolution） &
\textbf{Yes}（Theorem SE-1/2, Theorem 12.5） & \textbf{Yes}（M\_t
单调增长，no deletion） & \textbf{Arbitrary}（设计using 于physicalScience） \\
\end{longtable}

*AlphaZero at tabular caseunder has 策略迭代'sConvergenceProof；for 神via 网络情况as via 验性。

\subsubsection{4.2 Spring vs AlphaZero
Detailed Comparison}\label{spring-vs-alphazero-ux8be6ux7ec6ux5bf9ux6bd4}

\begin{longtable}[]{@{}
  >{\raggedright\arraybackslash}p{(\linewidth - 4\tabcolsep) * \real{0.2400}}
  >{\raggedright\arraybackslash}p{(\linewidth - 4\tabcolsep) * \real{0.4400}}
  >{\raggedright\arraybackslash}p{(\linewidth - 4\tabcolsep) * \real{0.3200}}@{}}
\toprule\noalign{}
\begin{minipage}[b]{\linewidth}\raggedright
Dimension
\end{minipage} & \begin{minipage}[b]{\linewidth}\raggedright
AlphaZero
\end{minipage} & \begin{minipage}[b]{\linewidth}\raggedright
Spring
\end{minipage} \\
\midrule\noalign{}
\endhead
\bottomrule\noalign{}
\endlastfoot
\textbf{Data Generation} & Synthetic (self-play)------environment fully known & Real (NEP
simulation)------physical规律not yet 知 \\
\textbf{Feedback} & Instant (game outcome) & Delayed (NEP batch training) \\
\textbf{Convergence Mechanism} & Policy iteration + MCTS & Lyapunov under 降 + Bayesian update \\
\textbf{Memory} & Replay buffer (with replacement) & 单调growing memory bank（non-删除） \\
\textbf{Convergence Guarantee} & 单调改进（tabular casealready Proof） & Finite-time fixed point (Theorem
SE-2, already Proof） \\
\textbf{Final State} & 原then on Optimal策略 & Self-consistent fixed point (possibly suboptimal) \\
\textbf{Exploration} & Dirichlet noise + visit count & Random exploration fraction ε +
annealed acceptance threshold \\
\end{longtable}

\subsubsection{4.3 Spring vs
Bayesian Optimization}\label{spring-vs-ux8d1dux53f6ux65afux4f18ux5316}

\begin{longtable}[]{@{}
  >{\raggedright\arraybackslash}p{(\linewidth - 4\tabcolsep) * \real{0.2400}}
  >{\raggedright\arraybackslash}p{(\linewidth - 4\tabcolsep) * \real{0.4400}}
  >{\raggedright\arraybackslash}p{(\linewidth - 4\tabcolsep) * \real{0.3200}}@{}}
\toprule\noalign{}
\begin{minipage}[b]{\linewidth}\raggedright
Dimension
\end{minipage} & \begin{minipage}[b]{\linewidth}\raggedright
Bayesian Optimization
\end{minipage} & \begin{minipage}[b]{\linewidth}\raggedright
Spring
\end{minipage} \\
\midrule\noalign{}
\endhead
\bottomrule\noalign{}
\endlastfoot
\textbf{Objective} & Optimize static black-box function & Evolve evaluation criteria + optimize model \\
\textbf{Surrogate Model} & Gaussian process & NEP student (neural network) \\
\textbf{Acquisition Function} & EI / UCB / Thompson & state valueFunction V(s) =
R̂(s)·ρ(s)·(1-C(s)) \\
\textbf{Granularity} & Pointwise & 逐Status（聚合多个Sample） \\
\textbf{Number of Objectives} & Single-objective (usually) &
Multi-objective (noise detection, expert disagreement, data diversity) \\
\textbf{Cost Structure} & Homogeneous & 异质（DFT Computation成本because Statuswhile 异） \\
\textbf{ConvergenceRate} & O(√(T·γ\_T))（Cumulative regret） & O(t\^{}\{-a\}) +
has 限时间终止Guarantee \\
\end{longtable}

\subsubsection{4.4 Spring vs
主动Learning}\label{spring-vs-ux4e3bux52a8ux5b66ux4e60}

\begin{longtable}[]{@{}
  >{\raggedright\arraybackslash}p{(\linewidth - 4\tabcolsep) * \real{0.2609}}
  >{\raggedright\arraybackslash}p{(\linewidth - 4\tabcolsep) * \real{0.3913}}
  >{\raggedright\arraybackslash}p{(\linewidth - 4\tabcolsep) * \real{0.3478}}@{}}
\toprule\noalign{}
\begin{minipage}[b]{\linewidth}\raggedright
Dimension
\end{minipage} & \begin{minipage}[b]{\linewidth}\raggedright
主动Learning
\end{minipage} & \begin{minipage}[b]{\linewidth}\raggedright
Spring
\end{minipage} \\
\midrule\noalign{}
\endhead
\bottomrule\noalign{}
\endlastfoot
\textbf{Query Strategy} & non-确定性 / QBC / Expected model change &
State Certification------at state-level聚合后决策 \\
\textbf{Model} & Single discriminative model & Multi-expert ensemble + student Model \\
\textbf{Query Decision} & Pointwise & Hierarchical------先state-level，again sample-level \\
\textbf{Noise Handling} & Oracle assumes clean (most) &
\textbf{主must 设计Objective}------explicitly detect label noise \\
\textbf{Label Complexity} & O(θ·d·log(1/ε)) & O(K\_S·(1/Δ\_min²)·log(1/ε)) \\
\textbf{Certification} & No & each 个lines动by state-levelguarantee certification \\
\end{longtable}

\begin{center}\rule{0.5\linewidth}{0.5pt}\end{center}

\subsection{5. Spring
Unique Contributions}\label{spring-ux72ecux6709ux7684ux8d21ux732e}

with under Yes Spring 完成while \textbf{no prior work has simultaneously achieved}:

\subsubsection{5.1
Self-Evolutionevaluation criteria（non-only 自改进策略）}\label{ux81eaux8fdbux5316ux8bc4ux4f30ux6807ux51c6ux4e0dux4ec5ux81eaux6539ux8fdbux7b56ux7565}

\textbf{Prior Work Limitations}: AlphaZero
自改进its 策略（such as 何under 棋），but evaluation criteria（游戏规then ）Yes固定's外部真理。主动Learning和Bayesian Optimizationhas Fixed Query Strategy/Acquisition Function。

\textbf{Spring's Innovation}: Gatekeeper \(S_t\)
\textbf{itself}YesLearningfor 象------it defines''什么Yes好'sData''。Systemnon-Yesat 固定evaluation criteriaunder 优化策略，while Yes\textbf{simultaneously evolvingevaluation criteria本身}。this is a二阶LearningProblem：Learningsuch as 何Learning判断Data质量。

\subsubsection{5.2
monotonically growing memory（non-替换重放）}\label{ux5355ux8c03ux589eux957fux8bb0ux5fc6ux4e0dux66ffux6362ux91cdux653e}

\textbf{Prior Work Limitations}:
which has 重放缓冲方案（AlphaZero、DQN、via 验重放）at 某种程度on 替换旧Memory------缓冲has 界，最旧/最non-has using 'sSampleis 丢弃。

\textbf{Spring's Innovation}:
\(M_t \subseteq M_{t+1}\)------Memoryonly 增长。no deletion策略意味着： 1.
No灾难性遗忘------早期Learning'sStructure永non-is 丢弃 2.
Memory单调性提供自by 'sstructural property（Lemma SE-1.2 requires must ） 3.
resurrection mechanism成as can be can ------分数Structures below threshold's旧Structure保留at memory bankin ，etc.待
Gatekeeper to mature and re-score them

\subsubsection{5.3 resurrection mechanism}\label{ux590dux6d3bux673aux5236}

\textbf{Prior Work Limitations}:
No前人Algorithm包含''复活'''s概念。is 拒绝'sSamplemust 么is 丢弃（主动Learning），must 么non-is 重访。

\textbf{Spring's Innovation}: ``Winter does not kill------it waits for spring''。Structures below threshold
\(\tau_{\text{keep}}\) 'sStructurenon-is 删除------它们进入dormancy。when  Gatekeeper
at 后续迭代in Evolution（\(S_{t+1}\)
can be can for 相同Structuregive 出更高分数）时，dormancyStructureis 复活并引入Training：

\begin{Shaded}
\begin{Highlighting}[]
\ControlFlowTok{for}\NormalTok{ (x, y, old\_score) }\KeywordTok{in}\NormalTok{ M\_t where old\_score }\OperatorTok{\textless{}}\NormalTok{ τ\_keep:}
\NormalTok{    new\_score }\OperatorTok{=}\NormalTok{ S\_\{t}\OperatorTok{+}\DecValTok{1}\NormalTok{\}(x, y)}
    \ControlFlowTok{if}\NormalTok{ new\_score }\OperatorTok{\textgreater{}=}\NormalTok{ τ\_keep:}
\NormalTok{        RESURRECT(x, y)  }\CommentTok{\# Spring resurrects it}
\end{Highlighting}
\end{Shaded}

this at 数学on 编码as ：Gatekeeper
Updatecan be with at 任何时间will Sample's评分推过阈值，变拒绝as 接受------this at 主动Learning（一旦拒绝i.e.永久丢弃）or Bayesian Optimization（一旦Evaluationi.e.消耗预算）in Yesnon-can be can 's。

\subsubsection{5.4 Coupled Dynamical System Analysis（Gatekeeper +
student co-evolution）}\label{ux8026ux5408ux52a8ux529bux7cfbux7edfux5206ux6790Gatekeeper-ux5b66ux751fux534fux540cux8fdbux5316}

\textbf{Prior Work Limitations}: Robbins-Monro analyzes a single recursion。AlphaZero
's策略迭代交替固定一个组件update the other（策略Evaluationand 策略改进之间NoCoupled Dynamical System Analysis）。

\textbf{Spring's Innovation}: Complete's耦合动力System形式化： - state space
\(\mathcal{Z} = \mathcal{F} \times \mathcal{M} \times \Theta\)（Product Space）
- update operator
\(\Phi = (\Phi_S, \Phi_M, \Phi_\theta)\)（because 果前馈Structure，block-triangular
Jacobian) - 双时间尺度 ODE Analysis（快慢方程组) - Lyapunov
Functionat Product Spaceon 's显式under 降 - \(\omega\)--limit set, basin of attraction, phase diagram

没has 前人Algorithmat 形式耦合动力System框架in Analysis两个Learning组件'sco-evolution。

\subsubsection{5.5 精确ConvergenceRateand  Polyak
平all }\label{ux7cbeux786eux6536ux655bux901fux7387ux4e0e-polyak-ux5e73ux5747}

\textbf{Prior Work Limitations}: Bayesian Optimizationgive 出cumulative regret bounds
\(O(\sqrt{T \cdot \gamma_T})\)
but non-give 出parameter convergence rates。主动Learninggive 出Label Complexitybut non-give 出LearningRate。

\textbf{Spring's Innovation}: - 强凸under student Rateas  \(O(t^{-a})\)（Theorem 11.1,
already Proof) - 收缩under  Gatekeeper Rateas  \(O(t^{-b})\)（Theorem 11.2,
already Proof) - Polyak-Ruppert Averagingwill Rate提升to  \(O(t^{-1})\)------matching
minimax lower bound \(\Omega(t^{-1})\) - 耦合Rateby student 瓶颈 \(O(t^{-a})\)
主导（Theorem 11.3) - has 限时间（non-渐近）界含显式常数（Theorem 11.6,
already Proof）

\subsubsection{5.6
形式化Selection Bias循环Analysis}\label{ux5f62ux5f0fux5316ux9009ux62e9ux504fux5deeux5faaux73afux5206ux6790}

\textbf{Prior Work Limitations}:
No前人Algorithm形式化Analysisits itselfDataSelection程序创建自强化偏差循环'sMechanism。

\textbf{Spring's Innovation}: Theorem 12.2
Proof------\textbf{Without reference set replay, Lyapunov
under 降at 渐近意义under Yesnon-can be can 's}------Gatekeeper
can reduce its apparent loss by only accepting''easy''Sample来降低its 表观损失，while non-真正改进。this Yesfor 任何self-selecting learning system's深刻限制，at 
Spring 之前not yet is 形式化Proof。

Specifically，Selection Bias循环：

\[\text{Gatekeeper } S_t \xrightarrow{\text{Selection}} \text{Memory } M_{t+1} \xrightarrow{\text{Training}} \text{NEP } f_{\theta_{t+1}} \xrightarrow{\text{Feedback}} \text{Gatekeeper } S_{t+1}\]

is 分解、量化并Proofas  Lyapunov under 降's必must 阻碍------then resolved through the combined mechanism of reference set replay +
重must 性采样's联合Mechanismis Resolution（Theorem 12.5, CompleteProof）。

\begin{center}\rule{0.5\linewidth}{0.5pt}\end{center}

\subsection{Summary}\label{ux603bux7ed3}

Spring
Self-evolution algorithmat 七个via 典数学Domain's交汇处运作：Stochastic Approximation、动力System、贝叶斯推断、at 线Learning、concentration inequalities、Information Theory和
minimax
Theory。its 核心架构创新------Two-Timescale Separation、参考集重放打破Selection Bias、monotonically growing memory、resurrection mechanism、Coupled Dynamical System Analysis------by this 些Domain'sCompleteProof链支撑。and which has 前人's决定性区别at 于：Spring
\textbf{simultaneously}Self-Evolutionits evaluation criteria和is EvaluationModel，formally proves convergence，并保持单调增长'sMemory（永non-遗忘）。

\begin{center}\rule{0.5\linewidth}{0.5pt}\end{center}

\emph{本谱系学by  SCX Theory组编制, 2026-06-28}

\end{document}
